% ============================================================
% Section 02: 量子力学简答题与概念题汇总
% ============================================================
\section{量子力学：简答题与概念题汇总}
\label{sec:qm-short-questions}

\begin{tipbox}{关键词标签}
叠加原理 · 定态 · 自旋 · 全同粒子 · 简并 · 束缚态 · 宇称 · 隧道效应 · 不确定关系 · 厄米算符 · 选择定则 · 微扰论 · 角动量耦合 · 占有数表象
\end{tipbox}

\begin{insightbox}{本章定位}
本章汇总量子力学各章节的核心概念与简答题，覆盖基本原理、算符理论、角动量与自旋、微扰论、全同粒子等主题。答案经过精炼，注重物理概念的准确表述与逻辑完整性。
\end{insightbox}

\vspace{0.5em}

\subsection{一、简答题}

\begin{enumerate}[label=\textbf{\arabic*.}, leftmargin=2em, itemsep=0.8em]

% --------------------------------------------------------
\item \textbf{试表述量子态的叠加原理并说明叠加系数是否依赖于时空变量及其理由。}

\textbf{叠加原理}：若 $\psi_1$ 和 $\psi_2$ 是体系可能的量子态，则它们的任意线性叠加
\[
    \psi = c_1\psi_1 + c_2\psi_2
\]
也是体系可能的状态，其中 $c_1, c_2$ 为复常数（一般地，$\psi = \sum_n c_n\psi_n$）。

\textbf{叠加系数不依赖于时空变量}。理由：叠加原理反映的是量子态所在希尔伯特空间的\textbf{线性结构}，系数 $c_n$ 描述的是各基态在叠加中所占的权重，是抽象的态矢量在给定基矢下的展开系数。若系数依赖于时空坐标，则叠加后的态将无法满足含时薛定谔方程 $i\hbar\,\partial\Psi/\partial t = \hat{H}\Psi$ 的线性齐次结构。

% --------------------------------------------------------
\item \textbf{一个量子体系处于定态的条件是什么？}

体系处于定态的充要条件为：
\begin{enumerate}[nosep, label=(\arabic*)]
    \item 体系的哈密顿量 $\hat{H}$ \textbf{不显含时间}；
    \item 体系波函数为 $\hat{H}$ 的某个\textbf{能量本征态}，即满足定态薛定谔方程 $\hat{H}\psi = E\psi$，从而含时波函数可分离变量为
    \[
        \Psi(\mathbf{r},t) = \psi(\mathbf{r})\,e^{-iEt/\hbar}.
    \]
\end{enumerate}
此时体系的一切力学量的概率分布及其期望值均不随时间改变。

% --------------------------------------------------------
\item \textbf{自旋可以在坐标表象中表示吗？}

\textbf{不可以}。自旋是粒子的\textbf{内禀属性}，没有经典对应，不是坐标和动量的函数，因此不能在坐标表象中写成 $\psi(\mathbf{r})$ 的形式。

自旋需在独立的\textbf{自旋空间}中用旋量（spinor）描述。自旋量子数为 $s$ 的粒子，其自旋态为 $(2s+1)$ 维列向量。例如电子（$s=1/2$）的自旋态为二分量旋量
\[
    \chi = \begin{pmatrix} \psi_\uparrow \\ \psi_\downarrow \end{pmatrix},
\]
其中 $|\psi_\uparrow|^2$ 和 $|\psi_\downarrow|^2$ 分别表示测得自旋投影为 $+\hbar/2$ 和 $-\hbar/2$ 的概率。

% --------------------------------------------------------
\item \textbf{描写全同粒子体系状态的波函数有何特点？}

全同粒子体系的波函数必须满足\textbf{交换对称性}：
\begin{itemize}[nosep]
    \item \textbf{费米子}（半整数自旋，$s=1/2,3/2,\dots$）：波函数对交换任意两个粒子是\textbf{反对称}的，即 $\hat{P}_{ij}\Psi = -\Psi$。这导致泡利不相容原理——两个费米子不能处于同一单粒子态。
    \item \textbf{玻色子}（整数自旋，$s=0,1,2,\dots$）：波函数对交换任意两个粒子是\textbf{对称}的，即 $\hat{P}_{ij}\Psi = +\Psi$。玻色子不受泡利原理限制，多个玻色子可占据同一量子态（玻色-爱因斯坦凝聚）。
\end{itemize}

% --------------------------------------------------------
\item \textbf{解释量子力学中的"简并"和"简并度"。}

\textbf{简并}：若力学量算符 $\hat{F}$ 的某个本征值 $f_n$ 对应\textbf{多个}线性独立的本征函数，则称该本征值是简并的。

\textbf{简并度}：属于同一个本征值的线性独立本征函数的个数，记为 $d_n$ 或 $f_n$。

\textbf{例子}：氢原子能级 $E_n$ 的简并度为 $n^2$（不考虑自旋），因为给定 $n$，$l$ 可取 $0,1,\dots,n-1$，每个 $l$ 对应 $2l+1$ 个 $m$ 值，总和为 $\sum_{l=0}^{n-1}(2l+1)=n^2$。

% --------------------------------------------------------
\item \textbf{分别说明什么样的状态是束缚态、简并态与负宇称态？}

\begin{itemize}[nosep]
    \item \textbf{束缚态}：粒子被势场限制在有限空间区域内运动，当 $r\to\infty$ 时波函数趋于零（$\psi\to 0$），能量本征值\textbf{分立}。如氢原子的束缚态、一维谐振子、无限深势阱中的粒子。
    \item \textbf{简并态}：多个线性独立的本征函数对应\textbf{同一个本征值}的态。如氢原子中 $(n,l,m)=(2,0,0)$ 与 $(2,1,0)$ 具有相同能量 $E_2$，是简并态。
    \item \textbf{负宇称态}：在空间反射（反演）变换 $\mathbf{r}\to -\mathbf{r}$ 下，波函数满足 $\hat{P}\psi(\mathbf{r}) = \psi(-\mathbf{r}) = -\psi(\mathbf{r})$，即宇称本征值为 $-1$。轨道角动量为奇数（$l=1,3,5,\dots$）的态具有负宇称。
\end{itemize}

% --------------------------------------------------------
\item \textbf{放射性指的是束缚在某些原子核中的更小粒子有一定的概率逃逸出来，你认为这与什么量子效应有关？}

与\textbf{量子隧道效应}（Quantum Tunneling Effect）有关。经典力学中，若粒子能量 $E$ 低于势垒高度 $V_0$，则粒子不可能穿透势垒。但量子力学中，粒子具有波动性，即使 $E<V_0$，仍有一定概率穿透势垒（或从势阱中逸出）。$\alpha$ 衰变中，$\alpha$ 粒子在核内形成后，通过隧穿核势垒逃逸，伽莫夫（Gamow）因子 $e^{-2G}$ 给出了穿透概率的量级估计。

% --------------------------------------------------------
\item \textbf{在量子力学中，能不能同时用粒子坐标和动量的确定值来描写粒子的量子状态？}

\textbf{不能}。坐标与动量是一对\textbf{共轭力学量}，满足海森堡不确定性原理
\[
    \Delta x \cdot \Delta p_x \ge \frac{\hbar}{2}.
\]
它们不能同时具有确定值。量子态可选用坐标表象 $\psi(x)$ 或动量表象 $\phi(p)$ 描述，但不能同时用两者的确定值来描写状态。这是量子力学与经典力学的根本区别之一。

% --------------------------------------------------------
\item \textbf{量子力学中的力学量算符有哪些性质？为什么需要这些性质？}

\textbf{核心性质}：
\begin{enumerate}[nosep, label=(\arabic*)]
    \item \textbf{线性性}：$\hat{F}(c_1\psi_1 + c_2\psi_2) = c_1\hat{F}\psi_1 + c_2\hat{F}\psi_2$；
    \item \textbf{厄米性}（自伴性）：$\hat{F}^\dagger = \hat{F}$，即对任意 $\psi,\phi$ 有 $\langle\psi|\hat{F}|\phi\rangle = \langle\phi|\hat{F}|\psi\rangle^*$。
\end{enumerate}

\textbf{必要性}：
\begin{itemize}[nosep]
    \item 线性性保证态叠加原理与算符作用的自洽性，使薛定谔方程保持线性；
    \item 厄米性保证本征值（即测量结果）为\textbf{实数}，且不同本征值对应的本征态\textbf{正交}，从而可构成希尔伯特空间的完备正交基。
\end{itemize}

% --------------------------------------------------------
\item \textbf{什么是塞曼效应？什么是斯达克效应？}

\begin{itemize}[nosep]
    \item \textbf{塞曼效应（Zeeman Effect）}：原子在外\textbf{磁场}中，由于磁矩与磁场的相互作用，能级发生分裂，导致光谱线分裂的现象。
    \begin{itemize}[nosep]
        \item \textbf{正常塞曼效应}：仅考虑轨道磁矩，自旋为零时，能级分裂为三条，分裂间距为 $\Delta E = \mu_B B$；
        \item \textbf{反常塞曼效应}：考虑电子自旋，分裂条数和间距更复杂，由朗德 $g$ 因子决定。
    \end{itemize}
    \item \textbf{斯达克效应（Stark Effect）}：原子在外\textbf{电场}中，能级发生分裂或位移的现象。氢原子的 $n=2$ 能级在电场中发生一级（线性）分裂，而基态 $n=1$ 只有二级（二次）斯达克效应。
\end{itemize}

% --------------------------------------------------------
\item \textbf{简述波函数的统计解释。}

波函数 $\Psi(\mathbf{r},t)$ 的模平方
\[
    \rho(\mathbf{r},t) = |\Psi(\mathbf{r},t)|^2
\]
表示粒子在时刻 $t$ 出现在空间位置 $\mathbf{r}$ 处的\textbf{概率密度}。具体地，$|\Psi(\mathbf{r},t)|^2\,d^3r$ 给出在体积元 $d^3r$ 内找到粒子的概率。

波函数必须满足\textbf{归一化条件} $\int|\Psi|^2\,d^3r = 1$，且可差一个任意整体相位因子 $e^{i\alpha}$（$\alpha$ 为实常数），因为概率密度不受相位影响。这是量子力学哥本哈根诠释的核心内容，由玻恩（Born）于1926年提出。

% --------------------------------------------------------
\item \textbf{对"轨道"和"电子云"的概念，量子力学的解释是什么？}

\begin{itemize}[nosep]
    \item \textbf{轨道}：旧量子论（玻尔模型）中电子绕核运动的经典概念。量子力学中\textbf{摒弃了轨道概念}，因为它与不确定性原理矛盾——若电子有确定的轨道，则同时具有确定的位置和动量，这违反了 $\Delta x\cdot\Delta p\ge \hbar/2$。
    \item \textbf{电子云}：量子力学中，电子在原子中的状态由波函数 $\psi_{nlm}(\mathbf{r})$ 描述，$|\psi_{nlm}(\mathbf{r})|^2$ 给出电子在空间各点出现的概率分布。将这一概率分布形象化即为\textbf{电子云}——云密处概率大，云疏处概率小。
\end{itemize}

% --------------------------------------------------------
\item \textbf{力学量在自身表象中的矩阵表示有何特点？}

若以力学量 $\hat{F}$ 的本征函数系 $\{|f_n\rangle\}$ 为基矢（即\textbf{自身表象}），则 $\hat{F}$ 的矩阵表示为\textbf{对角矩阵}，对角元即为其本征值：
\[
    F_{mn} = \langle f_m|\hat{F}|f_n\rangle = f_n\,\delta_{mn}.
\]
这是自身表象的定义性特征，也是表象理论中最简洁的情形。

% --------------------------------------------------------
\item \textbf{何为束缚态？}

束缚态是指粒子被势场约束在\textbf{有限区域}内运动的状态，其波函数在无穷远处趋于零：
\[
    \psi(\mathbf{r}) \xrightarrow{r\to\infty} 0,
\]
相应的能量本征值是\textbf{分立}的。典型的束缚态包括：无限深势阱中的粒子、谐振子、氢原子的负能态等。与之相对的是\textbf{散射态}（连续谱），其波函数在无穷远处不衰减，能量本征值连续。

% --------------------------------------------------------
\item \textbf{Stern--Gerlach 实验证实了什么？}

Stern--Gerlach 实验（1922年）证实了以下两点：
\begin{enumerate}[nosep, label=(\arabic*)]
    \item \textbf{空间量子化}：角动量在空间特定方向（外磁场方向）的投影是量子化的，只能取分立值；
    \item \textbf{电子自旋的存在}：银原子束在不均匀磁场中分裂为\textbf{两束}，说明电子具有内禀自旋角动量，且其投影只能取 $s_z = \pm\hbar/2$ 两个值。这与轨道角动量（分裂为奇数条）不同，直接证明了自旋 $s=1/2$ 的存在。
\end{enumerate}

% --------------------------------------------------------
\item \textbf{一个物理体系存在束缚态的条件是什么？}

存在束缚态的条件：势场必须能够"束缚"粒子，即至少存在一个能量本征值 $E$ 低于无穷远处的势能值 $V(\infty)$（通常取 $V(\infty)=0$ 为零点），使得粒子被局域在有限空间区域内。

\textbf{维数依赖性}：
\begin{itemize}[nosep]
    \item \textbf{一维}：对于任意吸引势 $V(x)<0$，无论多浅，都至少存在一个束缚态；
    \item \textbf{三维}：势阱必须足够深、足够宽，才能存在束缚态（存在临界深度条件）。
\end{itemize}

% --------------------------------------------------------
\item \textbf{两个对易的力学量是否一定同时确定？为什么？}

\textbf{不一定}。若 $[\hat{A},\hat{B}]=0$，则 $\hat{A}$ 和 $\hat{B}$ 具有\textbf{完备的共同本征函数系}，\textbf{原则上}可以同时被精确测定。但"同时确定"还取决于体系所处的\textbf{具体状态}：
\begin{itemize}[nosep]
    \item 若体系处于 $\hat{A}$ 和 $\hat{B}$ 的某个\textbf{共同本征态}，则两者同时具有确定值；
    \item 若体系处于叠加态（如 $\hat{A}$ 存在简并，态是若干简并本征态的叠加），则即使 $[\hat{A},\hat{B}]=0$，$\hat{B}$ 也可能不确定。
\end{itemize}
简言之：对易保证\textbf{可以}同时测定，不等于\textbf{一定}同时确定——后者取决于态的具体形式。

% --------------------------------------------------------
\item \textbf{什么是厄米矩阵，特点是什么？}

\textbf{厄米矩阵}（Hermitian Matrix）：满足 $F = F^\dagger$（即 $F^\dagger = (F^*)^T = F$）的方阵，矩阵元满足 $F_{mn} = F_{nm}^*$。

\textbf{特点}：
\begin{enumerate}[nosep, label=(\arabic*)]
    \item 对角元为实数：$F_{nn} = F_{nn}^*$；
    \item 本征值为\textbf{实数}；
    \item 不同本征值对应的本征矢\textbf{正交}；
    \item 可通过\textbf{幺正变换}对角化；
    \item 本征矢构成完备集，可用作希尔伯特空间的基。
\end{enumerate}

% --------------------------------------------------------
\item \textbf{何为选择定则？}

选择定则是指由初态 $|i\rangle$ 跃迁到末态 $|f\rangle$ 时，跃迁矩阵元 $\langle f|\hat{H}'|i\rangle \neq 0$ 所要求的\textbf{量子数条件}。它决定哪些跃迁是\textbf{允许}的、哪些是\textbf{禁戒}的。

\textbf{电偶极跃迁的选择定则}（原子）：
\[
    \Delta l = \pm 1, \qquad \Delta m = 0, \pm 1.
\]
例如氢原子 $2p\to 1s$ 是允许跃迁，而 $2s\to 1s$ 是禁戒的（电偶极近似下），需要通过双光子过程或高阶多极跃迁实现。

% --------------------------------------------------------
\item \textbf{在什么样的基组中，厄米算符是厄米矩阵？}

在\textbf{任何正交归一基组}中，厄米算符对应的矩阵都是厄米矩阵。即若基矢 $\{|n\rangle\}$ 满足 $\langle m|n\rangle = \delta_{mn}$，则矩阵元 $F_{mn} = \langle m|\hat{F}|n\rangle$ 自动满足 $F_{mn} = F_{nm}^*$。

这一性质是厄米算符定义的直接推论，与具体基组的选择无关——这正是狄拉克符号体系的优势：算符的厄米性是一个与表象无关的抽象性质。

% --------------------------------------------------------
\item \textbf{自旋如何得来？}

自旋的发现经历了以下阶段：
\begin{enumerate}[nosep, label=(\arabic*)]
    \item \textbf{实验驱动}：Stern--Gerlach 实验（1922）显示银原子束分裂为两束；碱金属光谱的\textbf{精细结构}和\textbf{反常塞曼效应}也无法用轨道角动量解释。
    \item \textbf{假设提出}：1925年，乌伦贝克（Uhlenbeck）和古兹密特（Goudsmit）提出电子具有内禀角动量——\textbf{自旋}，量子数 $s=1/2$，并引入自旋磁矩 $\boldsymbol{\mu}_s = -g_s\frac{e}{2m}\mathbf{S}$（$g_s\approx 2$）。
    \item \textbf{理论推导}：1928年，狄拉克（Dirac）从\textbf{相对论性量子力学方程}（狄拉克方程）中\textbf{自然导出}电子自旋，证明了自旋是相对论性量子效应的必然结果，而非人为假设。
\end{enumerate}

% --------------------------------------------------------
\item \textbf{我们知道平面单色波的电场能和磁场能相等，而在用微扰论计算发射系数和吸收系数时，我们为什么忽略了磁场对电子的作用？}

在光与原子相互作用的\textbf{电偶极近似}下：
\begin{enumerate}[nosep, label=(\arabic*)]
    \item 光波波长 $\lambda \sim 500\,$nm 远大于原子尺度 $a_0 \sim 0.05\,$nm，因此可将电磁波的位相因子近似为常数，相互作用主项为 $-e\mathbf{E}\cdot\mathbf{r}$；
    \item 磁场对电子的作用力为洛伦兹力 $-e\mathbf{v}\times\mathbf{B}$，其大小与电场作用力之比约为 $v/c \sim \alpha \approx 1/137$（精细结构常数），属于高阶小量；
    \item 因此，在最低阶近似（电偶极近似）下可忽略磁场作用，只保留电场项。
\end{enumerate}
磁偶极跃迁和电四极跃迁的强度比电偶极跃迁小 $(\alpha)^2 \sim 10^{-4}$ 量级。

% --------------------------------------------------------
\item \textbf{对于自旋为 $3/2$ 的粒子，其自旋本征函数应是几行一列的矩阵？}

自旋量子数 $s=3/2$ 时，$s_z$ 的本征值 $m_s$ 可取 $-3/2, -1/2, +1/2, +3/2$，共 $2s+1 = 4$ 个值。因此自旋本征函数为 \textbf{4 行 1 列}的旋量矩阵：
\[
    \chi_{m_s} = \begin{pmatrix} \chi_{+3/2} \\ \chi_{+1/2} \\ \chi_{-1/2} \\ \chi_{-3/2} \end{pmatrix},
    \quad
    \hat{S}_z = \frac{\hbar}{2}\begin{pmatrix} 3 & 0 & 0 & 0 \\ 0 & 1 & 0 & 0 \\ 0 & 0 & -1 & 0 \\ 0 & 0 & 0 & -3 \end{pmatrix}.
\]

% --------------------------------------------------------
\item \textbf{克莱布希--高登系数是为解决什么问题提出的？}

克莱布希--高登系数（Clebsch--Gordan Coefficients, C-G 系数）用于解决\textbf{两个角动量的耦合问题}。

当两个角动量 $\mathbf{j}_1$ 和 $\mathbf{j}_2$ 耦合成总角动量 $\mathbf{J} = \mathbf{j}_1 + \mathbf{j}_2$ 时，耦合表象基矢 $|JM\rangle$ 可通过 C-G 系数用无耦合表象基矢 $|j_1m_1\rangle|j_2m_2\rangle$ 展开：
\[
    |JM\rangle = \sum_{m_1,m_2} C_{j_1m_1,j_2m_2}^{JM}\, |j_1m_1\rangle\,|j_2m_2\rangle,
\]
其中 C-G 系数仅在 $m_1+m_2=M$ 时非零，并满足归一化和相位约定（Condon-Shortley 相位）。

% --------------------------------------------------------
\item \textbf{旧量子论存在哪些不足？}

\begin{enumerate}[nosep, label=(\arabic*)]
    \item \textbf{人为引入量子化条件}：如玻尔-索末菲量子化条件 $\oint p\,dq = nh$，缺乏理论自洽性；
    \item \textbf{无法解释光谱线强度}和选择定则；
    \item \textbf{对复杂系统失效}：无法处理多电子原子、分子体系；
    \item \textbf{无法解释波动性}：不能解释电子衍射、隧道效应等；
    \item \textbf{无法解释自旋和全同粒子统计性质}：泡利原理和玻色-爱因斯坦统计超出其框架；
    \item \textbf{保留经典轨道概念}：与不确定性原理矛盾。
\end{enumerate}

% --------------------------------------------------------
\item \textbf{简述变分法的思想。}

变分法的核心思想基于\textbf{瑞利-里茨原理}（Rayleigh-Ritz Principle）：对于任意归一化试探波函数 $\psi$，能量期望值满足
\[
    E[\psi] = \frac{\langle\psi|\hat{H}|\psi\rangle}{\langle\psi|\psi\rangle} \ge E_0,
\]
其中 $E_0$ 为基态能量。

\textbf{步骤}：
\begin{enumerate}[nosep, label=(\arabic*)]
    \item 选择含参变量的试探波函数 $\psi(\alpha_1, \alpha_2, \dots)$；
    \item 计算能量期望值 $E(\boldsymbol{\alpha})$；
    \item 对参数求极小值：$\partial E/\partial\alpha_i = 0$；
    \item 得到的 $E_{\min}$ 是基态能量的上界，对应的 $\psi$ 是基态的近似波函数。
\end{enumerate}
变分法在无法精确求解时非常有用，如氦原子基态、分子轨道计算等。

\end{enumerate}

\begin{enumerate}[label=\textbf{\arabic*.}, leftmargin=2em, itemsep=0.8em, resume]

% --------------------------------------------------------
\item \textbf{写出电子在 $S_z$ 表象下的三个泡利矩阵。}

电子自旋算符 $\hat{S}_i = \frac{\hbar}{2}\sigma_i$，其中 $\sigma_i$ 为泡利矩阵。在 $S_z$ 表象中：
\[
    \sigma_x = \begin{pmatrix} 0 & 1 \\ 1 & 0 \end{pmatrix},\quad
    \sigma_y = \begin{pmatrix} 0 & -i \\ i & 0 \end{pmatrix},\quad
    \sigma_z = \begin{pmatrix} 1 & 0 \\ 0 & -1 \end{pmatrix}.
\]

\textbf{基本性质}：$\sigma_i^2 = I$，$\sigma_i\sigma_j = i\varepsilon_{ijk}\sigma_k$（$i\neq j$），$\{\sigma_i,\sigma_j\} = 2\delta_{ij}I$。

% --------------------------------------------------------
\item \textbf{引入狄拉克符号的意义何在？}

狄拉克符号（bra-ket notation）的引入具有深远意义：
\begin{enumerate}[nosep, label=(\arabic*)]
    \item \textbf{表象无关}：$|\psi\rangle$ 表示抽象的态矢量，不依赖于具体坐标系或表象；
    \item \textbf{运算简洁}：内积 $\langle\phi|\psi\rangle$、矩阵元 $\langle\phi|\hat{F}|\psi\rangle$、外积 $|\psi\rangle\langle\phi|$ 等形式紧凑清晰；
    \item \textbf{方便表象变换}：如 $\psi(x) = \langle x|\psi\rangle$、$\phi(p) = \langle p|\psi\rangle$；
    \item \textbf{便于处理连续谱和分立谱的统一框架}：如 $\langle x|x'\rangle = \delta(x-x')$、$\langle n|m\rangle = \delta_{nm}$。
\end{enumerate}

% --------------------------------------------------------
\item \textbf{定态微扰论的适用范围是什么？}

定态微扰论适用于以下情况：
\begin{enumerate}[nosep, label=(\arabic*)]
    \item 哈密顿量可分解为 $\hat{H} = \hat{H}_0 + \hat{H}'$，其中 $\hat{H}_0$ 的本征值和本征态\textbf{已知或可解}；
    \item 微扰 $\hat{H}'$ "足够小"，满足
    \[
        \Bigl|\frac{\langle m^{(0)}|\hat{H}'|n^{(0)}\rangle}{E_n^{(0)} - E_m^{(0)}}\Bigr| \ll 1 \quad (m\neq n),
    \]
    即微扰矩阵元远小于能级间距；
    \item 对于\textbf{简并态}，需使用简并微扰论——先将 $\hat{H}'$ 在简并子空间中对角化；
    \item 适用于\textbf{分立谱}的定态问题，不直接处理连续谱或含时演化。
\end{enumerate}

% --------------------------------------------------------
\item \textbf{简述两个角动量耦合的三角形关系。}

两个角动量 $\mathbf{j}_1$ 与 $\mathbf{j}_2$ 耦合成总角动量 $\mathbf{J} = \mathbf{j}_1 + \mathbf{j}_2$ 时，总角动量量子数 $J$ 满足\textbf{三角形法则}：
\[
    |j_1 - j_2| \le J \le j_1 + j_2,
\]
且 $J$ 以步长 1 取值。总磁量子数 $M = m_1 + m_2$，取值范围为 $-J \le M \le J$。

\textbf{例子}：$j_1=1$，$j_2=1/2$ 时，$J$ 可取 $1/2$ 和 $3/2$。

% --------------------------------------------------------
\item \textbf{叙述量子力学的态叠加原理。}

若 $\psi_1, \psi_2, \dots, \psi_n$ 为体系可能的状态，则其任意复系数线性叠加
\[
    \psi = \sum_i c_i\psi_i
\]
也是体系可能的状态，其中 $c_i$ 为复常数。

\textbf{测量诠释}：若体系处于叠加态 $\psi$，测量某力学量时，体系将以概率 $|c_i|^2$（假设基矢已归一且正交）坍缩到某个本征态 $\psi_i$，测量结果对应该本征值。态叠加原理是量子力学中\textbf{相干性}和\textbf{概率性}的根源。

% --------------------------------------------------------
\item \textbf{能否由薛定谔方程直接导出自旋？简单塞曼效应是否可以证实自旋的存在？}

\textbf{不能}。标准的非相对论薛定谔方程
\[
    i\hbar\frac{\partial\Psi}{\partial t} = \Bigl[-\frac{\hbar^2}{2m}\nabla^2 + V(\mathbf{r})\Bigr]\Psi
\]
不包含自旋自由度。自旋需作为额外假设引入，或从\textbf{狄拉克方程}（相对论性量子力学）中自然导出。

\textbf{简单（正常）塞曼效应不能证实自旋}，因为它仅考虑轨道磁矩与磁场的耦合，分裂为三条等间距谱线，可用 $l$ 的量子化解释。\textbf{反常塞曼效应}（分裂条数和间距不符合正常塞曼效应）以及 Stern--Gerlach 实验才直接证实了自旋的存在。

% --------------------------------------------------------
\item \textbf{非简并定态微扰论的计算公式是什么？写出其适用条件。}

设 $\hat{H}_0|n^{(0)}\rangle = E_n^{(0)}|n^{(0)}\rangle$ 已知，$\hat{H} = \hat{H}_0 + \hat{H}'$。

\textbf{一级能量修正}：
\[
    E_n^{(1)} = \langle n^{(0)}|\hat{H}'|n^{(0)}\rangle = H'_{nn}
\]

\textbf{一级波函数修正}：
\[
    |n^{(1)}\rangle = \sum_{m\neq n} \frac{H'_{mn}}{E_n^{(0)} - E_m^{(0)}}\,|m^{(0)}\rangle
\]

\textbf{二级能量修正}：
\[
    E_n^{(2)} = \sum_{m\neq n} \frac{|H'_{mn}|^2}{E_n^{(0)} - E_m^{(0)}}
\]

\textbf{适用条件}：$|H'_{mn}/(E_n^{(0)}-E_m^{(0)})| \ll 1$（$m\neq n$），且 $E_n^{(0)}$ 为非简并能级。

% --------------------------------------------------------
\item \textbf{波函数的量纲是否与表象有关？举例说明。}

\textbf{有关}。波函数量纲由归一化条件决定：

\textbf{坐标表象}：$\int|\psi(x)|^2\,dx = 1$（一维），故 $[\psi(x)] = L^{-1/2}$；三维情况下 $[\psi(\mathbf{r})] = L^{-3/2}$。

\textbf{动量表象}：$\int|\phi(p)|^2\,dp = 1$，因 $[p] = MLT^{-1}$，故 $[\phi(p)] = (MLT^{-1})^{-1/2}$。

\textbf{能量表象}（分立谱）：$\sum_n|c_n|^2 = 1$，展开系数 $c_n = \langle n|\psi\rangle$ 为无量纲纯数。

因此，波函数（或展开系数）的量纲依赖于所选表象的变量类型和维度。

% --------------------------------------------------------
\item \textbf{对线性谐振子定态问题，旧量子论与量子力学的结论存在哪些根本区别？}

\begin{center}
\renewcommand{\arraystretch}{1.4}
\begin{tabular}{|l|c|c|}
\hline
\textbf{性质} & \textbf{旧量子论} & \textbf{量子力学} \\
\hline
能量本征值 & $E_n = n\hbar\omega$ ($n=0,1,2,\dots$) & $E_n = (n+\frac{1}{2})\hbar\omega$ \\
零点能 & $E_0 = 0$ & $E_0 = \frac{1}{2}\hbar\omega \neq 0$ \\
粒子位置 & 有确定轨道（经典椭圆） & 无轨道，只有概率分布 \\
经典禁区 & 粒子不可能进入 & 有一定概率隧穿进入 \\
\hline
\end{tabular}
\end{center}

零点能的存在是量子力学的重要预言，已被实验（如分子光谱、低温比热）证实。

% --------------------------------------------------------
\item \textbf{简述氢原子的一级斯塔克效应。}

氢原子在外电场 $\mathcal{E}$ 中，微扰哈密顿量为 $\hat{H}' = e\mathcal{E}z = e\mathcal{E}r\cos\theta$（设电场沿 $z$ 轴）。

\textbf{$n=2$ 能级}（四重简并：$2s, 2p_0, 2p_{+1}, 2p_{-1}$）发生\textbf{一级（线性）斯塔克效应}：
\begin{itemize}[nosep]
    \item 在简并子空间中，$\hat{H}'$ 的矩阵元仅 $ \langle 2s|z|2p_0\rangle$ 和 $ \langle 2p_0|z|2s\rangle$ 非零；
    \item 对角化后，$2s$ 与 $2p_0$ 耦合成两个新态，能级分裂为
    \[
        \Delta E = \pm 3e\mathcal{E}a_0,
    \]
    其中 $a_0$ 为玻尔半径；
    \item $2p_{\pm 1}$ 态在一级微扰下能量不变（仍为 $E_2$）。
\end{itemize}

\textbf{基态 $n=1$}：只有 $1s$ 态，$\langle 1s|z|1s\rangle = 0$（球对称），故\textbf{无一级斯塔克效应}，只有二级（二次）效应 $\Delta E \propto \mathcal{E}^2$。

% --------------------------------------------------------
\item \textbf{什么是耦合表象？}

两个角动量 $\hat{\mathbf{J}}_1$ 和 $\hat{\mathbf{J}}_2$ 耦合成总角动量 $\hat{\mathbf{J}} = \hat{\mathbf{J}}_1 + \hat{\mathbf{J}}_2$。以 $\hat{\mathbf{J}}^2$ 和 $\hat{J}_z$ 的共同本征态 $|JM\rangle$ 为基矢的表象称为\textbf{耦合表象}（coupled representation）。

与之相对的是\textbf{无耦合表象}（uncoupled representation），其基矢为 $|j_1m_1\rangle|j_2m_2\rangle$（即 $\hat{\mathbf{J}}_1^2, \hat{J}_{1z}, \hat{\mathbf{J}}_2^2, \hat{J}_{2z}$ 的共同本征态）。两种表象通过 C-G 系数联系。

% --------------------------------------------------------
\item \textbf{给出线性谐振子定态波函数的递推公式。}

利用产生算符 $a^\dagger$ 和湮灭算符 $a$：
\[
    \psi_{n+1}(x) = \frac{1}{\sqrt{n+1}}\,a^\dagger\,\psi_n(x),\qquad
    \psi_{n-1}(x) = \frac{1}{\sqrt{n}}\,a\,\psi_n(x).
\]

显式的递推关系为
\[
    \psi_{n+1}(x) = \sqrt{\frac{2}{n+1}}\,\alpha x\,\psi_n(x) - \sqrt{\frac{n}{n+1}}\,\psi_{n-1}(x),
\]
其中 $\alpha = \sqrt{m\omega/\hbar}$。

基态波函数：$\psi_0(x) = (\alpha/\sqrt{\pi})^{1/2}\,e^{-\alpha^2x^2/2}$。

% --------------------------------------------------------
\item \textbf{对于氢原子，其偶极跃迁的选择定则对主量子数 $n$ 是否存在限制？为什么？}

\textbf{不存在限制}。电偶极跃迁的选择定则仅对角量子数 $l$ 和磁量子数 $m$ 有要求：
\[
    \Delta l = \pm 1, \qquad \Delta m = 0, \pm 1.
\]
主量子数 $n$ 不受限制，只要跃迁矩阵元 $\langle n'l'm'|e\mathbf{r}|nlm\rangle \neq 0$ 且满足上述 $\Delta l, \Delta m$ 条件，跃迁即可发生。例如 $3p\to 1s$、$4d\to 2p$、$5f\to 3d$ 等都是允许的跃迁。

% --------------------------------------------------------
\item \textbf{何谓力学量完全集？}

一组彼此\textbf{对易}且\textbf{函数独立}的力学量算符集合 $\{\hat{A}, \hat{B}, \hat{C}, \dots\}$，若它们的共同本征态足以\textbf{唯一确定}体系的状态（即消除所有简并），则称该集合为\textbf{力学量完全集}（Complete Set of Commuting Observables, CSCO）。

\textbf{例子}：氢原子可用 $\{\hat{H}, \hat{\mathbf{L}}^2, \hat{L}_z\}$ 作为力学量完全集，其共同本征函数为 $\psi_{nlm}(r,\theta,\varphi)$，用量子数 $(n,l,m)$ 完全标记一个状态。考虑自旋后，完全集扩展为 $\{\hat{H}, \hat{\mathbf{L}}^2, \hat{L}_z, \hat{S}_z\}$。

% --------------------------------------------------------
\item \textbf{有关角动量的定义，我们学过哪两种？哪一种更广泛？自旋角动量是按哪一种定义的？}

\begin{enumerate}[nosep, label=(\arabic*)]
    \item \textbf{轨道角动量}：$\hat{\mathbf{L}} = \hat{\mathbf{r}} \times \hat{\mathbf{p}}$，由经典对应通过正则量子化得到，量子数 $l$ 取整数 $0,1,2,\dots$。
    \item \textbf{自旋角动量}：内禀角动量，无经典对应，量子数 $s$ 可为整数或半整数 $0, 1/2, 1, 3/2, \dots$。
\end{enumerate}

\textbf{更广泛的定义}：通过\textbf{对易关系}定义的角动量
\[
    [\hat{J}_i, \hat{J}_j] = i\hbar\varepsilon_{ijk}\hat{J}_k,
\]
这是角动量最普遍的定义。轨道角动量和自旋角动量都满足此对易关系，但自旋角动量包含了半整数情况，因此更为广泛。自旋角动量就是按此对易关系定义的。

% --------------------------------------------------------
\item \textbf{说明在定态问题中，定态能量的最小值不可能低于势能的最低值。}

利用能量期望值分解：
\[
    E[\psi] = \langle\psi|\hat{H}|\psi\rangle = \langle\hat{T}\rangle + \langle V\rangle.
\]
由于动能算符 $\hat{T} = -\frac{\hbar^2}{2m}\nabla^2$ 是正定的，对任意态都有 $\langle\hat{T}\rangle \ge 0$。因此
\[
    E[\psi] \ge \langle V\rangle \ge V_{\min}.
\]
对定态而言，$E$ 为能量本征值，故 $E \ge V_{\min}$。定态能量的最小值不可能低于势能的最低值。

% --------------------------------------------------------
\item \textbf{简述占有数表象。}

占有数表象（Number Representation）将谐振子的激发量子视为"粒子"（声子或光子），以粒子数算符 $\hat{N} = a^\dagger a$ 的本征态 $|n\rangle$ 作为基矢：
\[
    \hat{N}|n\rangle = n|n\rangle, \quad n = 0,1,2,\dots
\]

产生和湮灭算符的作用：
\[
    a^\dagger|n\rangle = \sqrt{n+1}\,|n+1\rangle,\qquad
    a|n\rangle = \sqrt{n}\,|n-1\rangle.
\]

\textbf{意义}：在二次量子化（多体量子理论）中，占有数表象是描述全同粒子系统的基本语言，直接体现玻色子/费米子的统计性质。

% --------------------------------------------------------
\item \textbf{何为偶极近似？}

原子与电磁波相互作用时，若光的波长 $\lambda$ 远大于原子尺度 $a_0$，在原子范围内电磁波的位相变化可忽略，即位相因子 $e^{i\mathbf{k}\cdot\mathbf{r}} \approx 1$。此时相互作用哈密顿量近似为
\[
    \hat{H}' = -e\mathbf{E}\cdot\mathbf{r} = -\mathbf{d}\cdot\mathbf{E},
\]
其中 $\mathbf{d} = e\mathbf{r}$ 为电偶极矩。这种近似称为\textbf{偶极近似}（Dipole Approximation）或\textbf{长波近似}。

适用条件：$ka_0 \ll 1$，即 $\lambda \gg a_0$（对可见光，$\lambda \sim 10^3 a_0$，条件极好满足）。

% --------------------------------------------------------
\item \textbf{一个物理体系存在束缚态的条件是什么？}

（同第16题，参见上文。）核心条件：势阱能容纳至少一个能量低于无穷远处势能值的分立本征态。一维中任意吸引势至少有一个束缚态；三维中需势阱足够深宽。

% --------------------------------------------------------
\item \textbf{对于旧量子论中氢原子的"轨道"，量子力学的解释是什么？}

（同第12题，参见上文。）旧量子论的"轨道"被量子力学中的\textbf{概率密度分布}取代。电子状态由波函数 $\psi_{nlm}(\mathbf{r})$ 描述，$|\psi_{nlm}|^2$ 给出空间概率分布（电子云）。$|\psi|^2$ 的角向部分形成\textbf{电子云}的"形状"（s轨道球对称、p轨道哑铃形等），径向部分决定电子到核的距离分布。

\end{enumerate}

\vspace{1em}
\subsection{二、计算题（思路与解答）}

\begin{enumerate}[label=\textbf{\arabic*.}, leftmargin=2em, itemsep=0.8em]

% --------------------------------------------------------
\item \textbf{粒子的定态能量和波函数计算（以一维无限深势阱为例）}

\textbf{模型}：
\[
    V(x) = \begin{cases} 0, & 0 < x < a \\ \infty, & x \le 0 \text{ 或 } x \ge a \end{cases}
\]

\textbf{求解}：
\begin{enumerate}[nosep, label=Step \arabic*.]
    \item \textbf{定态薛定谔方程}（阱内 $V=0$）：
    \[
        -\frac{\hbar^2}{2m}\frac{d^2\psi}{dx^2} = E\psi \quad\Longrightarrow\quad
        \frac{d^2\psi}{dx^2} + k^2\psi = 0, \quad k = \frac{\sqrt{2mE}}{\hbar}.
    \]
    \item \textbf{通解}：$\psi(x) = A\sin kx + B\cos kx$。
    \item \textbf{边界条件}：$\psi(0) = 0$ 得 $B=0$；$\psi(a) = 0$ 得 $ka = n\pi$（$n=1,2,3,\dots$），$n=0$ 给出平庸解，舍去。
    \item \textbf{归一化}：$\int_0^a |\psi|^2\,dx = 1$ 得 $A = \sqrt{2/a}$。
\end{enumerate}

\textbf{结果}：
\[
    \boxed{\psi_n(x) = \sqrt{\frac{2}{a}}\sin\Bigl(\frac{n\pi x}{a}\Bigr)}, \qquad
    \boxed{E_n = \frac{n^2\pi^2\hbar^2}{2ma^2}}, \quad n=1,2,3,\dots
\]

\textbf{物理意义}：能量量子化，基态 $n=1$ 有零点能 $E_1 = \frac{\pi^2\hbar^2}{2ma^2}$；波函数在阱内有 $n-1$ 个节点。

% --------------------------------------------------------
\item \textbf{厄米算符及其矩阵的计算}

\textbf{例题}：验证矩阵 $A = \begin{pmatrix} 1 & i \\ -i & 2 \end{pmatrix}$ 的厄米性并求本征值。

\textbf{解}：
\begin{enumerate}[nosep, label=Step \arabic*.]
    \item \textbf{验证厄米性}：
    \[
        A^\dagger = (A^*)^T = \begin{pmatrix} 1 & i \\ -i & 2 \end{pmatrix}^T = \begin{pmatrix} 1 & -i \\ i & 2 \end{pmatrix}^T \,?\,
    \]
    更正：直接计算共轭转置
    \[
        A^\dagger = \begin{pmatrix} 1^* & (-i)^* \\ i^* & 2^* \end{pmatrix}^T
        = \begin{pmatrix} 1 & i \\ -i & 2 \end{pmatrix} = A.
    \]
    故 $A$ 为厄米矩阵。
    \item \textbf{久期方程} $\det(A - \lambda I) = 0$：
    \[
        \begin{vmatrix} 1-\lambda & i \\ -i & 2-\lambda \end{vmatrix}
        = (1-\lambda)(2-\lambda) - (i)(-i)
        = (1-\lambda)(2-\lambda) - 1 = 0.
    \]
    展开：$\lambda^2 - 3\lambda + 1 = 0$。
    \item \textbf{本征值}：
    \[
        \lambda_{1,2} = \frac{3 \pm \sqrt{5}}{2}.
    \]
    均为实数，符合厄米矩阵性质。
\end{enumerate}

% --------------------------------------------------------
\item \textbf{自旋投影的本征方程及其概率分布的计算}

自旋沿方向 $\mathbf{n} = (\sin\theta\cos\phi, \sin\theta\sin\phi, \cos\theta)$ 的投影算符为
\[
    \hat{S}_n = \hat{\mathbf{S}}\cdot\mathbf{n} = \frac{\hbar}{2}(\sigma_x n_x + \sigma_y n_y + \sigma_z n_z).
\]

对于电子（$s=1/2$）：
\[
    S_n = \frac{\hbar}{2}\begin{pmatrix} \cos\theta & \sin\theta\,e^{-i\phi} \\ \sin\theta\,e^{i\phi} & -\cos\theta \end{pmatrix}.
\]

\textbf{本征方程} $S_n\chi = \lambda\chi$ 给出本征值 $\lambda = \pm\hbar/2$。

\textbf{本征旋量}：
\[
    \chi_+ = \begin{pmatrix} \cos\frac{\theta}{2} \\ \sin\frac{\theta}{2}\,e^{i\phi} \end{pmatrix},\qquad
    \chi_- = \begin{pmatrix} -\sin\frac{\theta}{2}\,e^{-i\phi} \\ \cos\frac{\theta}{2} \end{pmatrix}.
\]

\textbf{概率分布}：若体系处于 $|\chi_+\rangle$ 态，测量 $S_z$ 得到 $+\hbar/2$ 的概率为
\[
    P\bigl(S_z = +\tfrac{\hbar}{2}\bigr) = |\langle\uparrow|\chi_+\rangle|^2 = \cos^2\frac{\theta}{2},
\]
得到 $-\hbar/2$ 的概率为
\[
    P\bigl(S_z = -\tfrac{\hbar}{2}\bigr) = |\langle\downarrow|\chi_+\rangle|^2 = \sin^2\frac{\theta}{2}.
\]

\end{enumerate}
